\documentclass[11pt]{article}

\usepackage[margin=1in]{geometry}
\usepackage[T1]{fontenc}
\usepackage[utf8]{inputenc}
\usepackage{lmodern}
\usepackage{microtype}
\usepackage{amsmath,amssymb,amsthm,mathtools}
\usepackage{booktabs}
\usepackage{enumitem}
\usepackage{xurl}
\usepackage[hidelinks]{hyperref}

\newtheorem{definition}{Definition}
\newtheorem{lemma}{Lemma}
\newtheorem{theorem}{Theorem}
\newtheorem{corollary}{Corollary}
\newtheorem{remark}{Remark}

\newcommand{\TV}{\mathrm{TV}}
\newcommand{\Dinf}{D_{\infty}}
\newcommand{\ML}{\mathcal{L}_{\max}}

\title{MC-EQ Under Support Discipline:\\
Transcript-TV Bounds, Reference-Ratio Guarantees, and Off-Support Tag Attacks}
\author{}
\date{Draft v0.54 (June 18, 2026; archive rev0897)}

\begin{document}
\maketitle

\begin{abstract}
A history-conditioned transition law can be close in total variation to a public destination-independent reference process while still placing a small amount of probability on destination-specific outcomes that the reference declares impossible.  That distinction is decisive.  Local total-variation equalization composes to a valid transcript-TV bound, but it does not by itself give a finite likelihood-ratio, pure differential-privacy, or maximal-leakage certificate.

This note repairs the MC-EQ interface by separating two modes.  \emph{TV mode} needs no support assumption and yields only reference-chain and pairwise transcript-TV guarantees.  \emph{Reference-ratio mode} additionally requires zero off-support mass and a positive reference floor; only then do likelihood-ratio and maximal-leakage bounds follow.  An explicit $N$-destination tag attack satisfies one-step MC-EQ with slack $\delta$, has pairwise TV only $\delta$, yet has infinite $D_\infty$ to the reference and maximal leakage $\log_2(1+(N-1)\delta)$.  For $N=256$ and $\delta=0.02$, that is $2.608809$ bits in one step.

The archive audit also withdraws a concrete packet binding: the maintained worked-example value $0.056345507023$ is a sum of three observation-attenuation terms, not an MC-EQ calculation.  No deployment MC-EQ certificate is shipped in this revision.
\end{abstract}

\paragraph{Scope and units.}
All logarithms in endpoint budgets are base two.  The protected object is a declared observer transcript, including every visible stop, abort, retry, and side-information symbol.  A theorem for a state-only projection is not a theorem for a richer network transcript.

\paragraph{Artifact posture.}
This revision ships theorem text and executable hostile-review arithmetic in the archive checker.  It does not ship a deployment kernel estimate, trace-derived slack certificate, or source-bound MC-EQ release card.

\section{Model and the two MC-EQ modes}
Let $D$ be a finite destination family.  For destination $d\in D$, an observer sees a finite transcript
\[
Z_{0:T}=(Z_0,Z_1,\ldots,Z_T).
\]
The alphabet must include any observer-visible termination or abort outcome.  Write $h_t=z_{0:t}$ for the visible history.  A public destination-independent reference process consists of an initial law $\mu$ and conditional kernels $K_t(\cdot\mid h_t)$.  The Markov specialization is $K_t(\cdot\mid h_t)=K(\cdot\mid x_t)$ for a public state $x_t$ carried by the transcript.

\begin{definition}[TV-MC-EQ]
The mechanism satisfies TV-MC-EQ with slacks $\varepsilon_0,\ldots,\varepsilon_{T-1}$ when, for every destination and every reachable visible history,
\[
\TV\!\left(P_d(Z_{t+1}\in\cdot\mid h_t),K_t(\cdot\mid h_t)\right)\le \varepsilon_t.
\]
No support-containment statement is part of this definition.
\end{definition}

\begin{definition}[Support-disciplined MC-EQ]
TV-MC-EQ is \emph{support disciplined} when, for every $d,t,h_t$,
\[
P_d(y\mid h_t)=0\quad\text{whenever}\quad K_t(y\mid h_t)=0,
\]
and every positive reference atom is bounded below by a declared number $q_t>0$:
\[
K_t(y\mid h_t)>0\quad\Longrightarrow\quad K_t(y\mid h_t)\ge q_t.
\]
\end{definition}

The first mode supports TV statements.  The second adds the absolute-continuity and denominator evidence required for ratio statements.  Conflating them was the central error in the earlier draft.

\section{What TV-MC-EQ actually proves}
Let $\mathcal L_d$ be the law of $Z_{0:T}$ under destination $d$, and let $\mathcal R$ be the reference transcript law.  Let $\gamma_d=\TV(P_d(Z_0\in\cdot),\mu)$.

\begin{theorem}[History-robust transcript-TV composition]
Define
\[
r_d=1-(1-\gamma_d)\prod_{t=0}^{T-1}(1-\varepsilon_t).
\]
Then
\[
\TV(\mathcal L_d,\mathcal R)\le r_d
\le \gamma_d+\sum_{t=0}^{T-1}\varepsilon_t.
\]
Consequently, for destinations $d,d'$,
\[
\TV(\mathcal L_d,\mathcal L_{d'})
\le \min\{1,r_d+r_{d'}\}.
\]
\end{theorem}

\begin{proof}
Maximally couple the initial laws, which match with probability at least $1-\gamma_d$.  Conditional on a matched visible history, maximally couple the next mechanism output to the reference kernel; the conditional mismatch probability is at most $\varepsilon_t$.  Therefore the probability that the entire transcript remains matched is at least $(1-\gamma_d)\prod_t(1-\varepsilon_t)$.  The coupling characterization of TV gives the first bound.  The additive envelope is the union bound, and the cross-destination result is triangle inequality through $\mathcal R$.
\end{proof}

\begin{corollary}[Binary distinguishing]
For equal priors on $d,d'$, the optimal destination-classification success probability from this transcript is at most
\[
\frac{1+\TV(\mathcal L_d,\mathcal L_{d'})}{2}.
\]
Thus the advantage over random guessing is one half of the pairwise-TV bound, not the full TV value.
\end{corollary}

This theorem remains valid with destination-specific outcomes outside the reference support.  That is precisely why it cannot be silently upgraded to a likelihood-ratio theorem.

\section{Off-support tags: small TV, infinite reference ratio}
\begin{theorem}[Alphabet-amplified off-support tag attack]
Fix $N\ge2$ destinations and $0<\delta<1$.  Let the reference output be the single core symbol $c$, so $Q(c)=1$.  For each destination $d\in\{1,\ldots,N\}$, introduce a unique tag $\tau_d$ and define
\[
P_d(c)=1-\delta,\qquad P_d(\tau_d)=\delta.
\]
Then all of the following hold:
\begin{enumerate}[leftmargin=1.5em]
\item $\TV(P_d,Q)=\delta$ for every destination, so one-step TV-MC-EQ holds;
\item $\TV(P_d,P_{d'})=\delta$ for every $d\ne d'$;
\item $\Dinf(P_d\|Q)=\infty$ because $Q(\tau_d)=0<P_d(\tau_d)$;
\item for finite alphabets, maximal leakage satisfies
\[
\ML(D\!\to\! Z)=\log_2\sum_z\max_d P_d(z)
=\log_2\!\bigl(1+(N-1)\delta\bigr);
\]
\item under a uniform destination prior, Bayes recovery is
\[
\Pr[\widehat D=D]=\delta+\frac{1-\delta}{N},
\]
which is a multiplicative gain $1+(N-1)\delta$ over the prior success $1/N$.
\end{enumerate}
\end{theorem}

\begin{proof}
The TV statements follow by direct $\ell_1$ calculation.  The destination-specific tag has positive mechanism mass and zero reference mass, giving infinite max-divergence.  The discrete maximal-leakage identity is the closed form of Issa, Kamath, and Wagner \cite{issa_maxleak}; the core contributes $1-\delta$ to the sum of column maxima and the $N$ tags contribute $N\delta$.  A tag reveals the destination exactly, while the core symbol leaves the posterior uniform.
\end{proof}

For the pinned drill $N=256$ and $\delta=0.02$,
\[
\ML(D\!\to\! Z)=\log_2(6.1)=2.608809242676\ \text{bits},
\]
and uniform-prior recovery is $0.023828125$ instead of $1/256=0.00390625$.  Pairwise TV remains only $0.02$.  This is not a contradiction: pairwise testing and many-destination maximal leakage are different claims.

For $T$ conditionally independent repetitions with the same hidden destination, the exact tag-attack values are
\[
\ML(D\!\to\! Z_{1:T})=
\log_2\!\left(N-(N-1)(1-\delta)^T\right)
\]
and
\[
\Pr[\widehat D=D]=1-(1-\delta)^T+\frac{(1-\delta)^T}{N}.
\]
At $T=6$, the pinned attack reaches $4.912180044586$ bits and recovery $0.117617940936$.

\section{The repaired reference-ratio theorem}
The support condition is not decorative.  It turns the TV statement into a finite pointwise statement.

\begin{lemma}[The correct singleton bound]
For discrete $P,Q$,
\[
\TV(P,Q)\le\delta\quad\Longrightarrow\quad |P(y)-Q(y)|\le\delta
\]
for every atom $y$.
\end{lemma}

\begin{proof}
Total variation is the supremum of $|P(A)-Q(A)|$ over events $A$; take $A=\{y\}$.  The earlier factor-$2$ bound obtained from the full $\ell_1$ norm was valid but unnecessarily loose.
\end{proof}

\begin{theorem}[Support-disciplined maximal leakage and pairwise ratio]
Assume support-disciplined MC-EQ.  Suppose also that the initial laws obey a common-reference ratio
\[
P_d(Z_0=z)\le 2^{b_0}\mu(z)
\quad\text{for every }d,z.
\]
Then
\[
\ML(D\!\to\! Z_{0:T})
\le b_0+\sum_{t=0}^{T-1}\log_2\!\left(1+\frac{\varepsilon_t}{q_t}\right).
\]

If the initial laws are identical and $\varepsilon_t<q_t$ for every $t$, then every pair $d,d'$ also satisfies
\[
\Dinf(\mathcal L_d\|\mathcal L_{d'})
\le
\sum_{t=0}^{T-1}
\log_2\!\left(\frac{q_t+\varepsilon_t}{q_t-\varepsilon_t}\right).
\]
\end{theorem}

\begin{proof}
The singleton lemma and the floor give
\[
P_d(y\mid h_t)\le K_t(y\mid h_t)+\varepsilon_t
\le \left(1+\frac{\varepsilon_t}{q_t}\right)K_t(y\mid h_t).
\]
Multiplying these history-conditioned inequalities along a transcript gives a common reference-path ratio.  Summing the maximum conditional path probability over transcripts yields the maximal-leakage bound.  For the pairwise statement,
\[
P_d(y\mid h_t)\le K_t(y\mid h_t)+\varepsilon_t,
\qquad
P_{d'}(y\mid h_t)\ge K_t(y\mid h_t)-\varepsilon_t,
\]
and the worst ratio occurs at the smallest reference mass $q_t$.  Multiplication over the visible history gives the transcript max-divergence bound.  No conditional-independence assumption is used.
\end{proof}

\begin{corollary}[Uniform public covers]
If every reference row is uniform on at most $k_t$ outcomes, then $q_t\ge1/k_t$.  Hence
\[
\ML(D\!\to\! Z_{0:T})
\le b_0+\sum_t\log_2(1+k_t\varepsilon_t),
\]
and, when $k_t\varepsilon_t<1$,
\[
\Dinf(\mathcal L_d\|\mathcal L_{d'})
\le\sum_t\log_2\!\left(\frac{1+k_t\varepsilon_t}{1-k_t\varepsilon_t}\right).
\]
\end{corollary}

A six-step example with $k=8$ and $\varepsilon=0.005$ gives
\[
\ML\le 6\log_2(1.04)=0.339501170198\ \text{bits}
\]
and a pairwise max-divergence bound of
\[
6\log_2(1.04/0.96)=0.692863304520\ \text{bits}
\]
when written in the normalized $1\pm k\varepsilon$ form.

\section{The worked-example binding that is withdrawn}
The maintained worked example contains a line item named
\texttt{fallback\_tiered\_vector} with public value
\[
0.056345507023
=0.046533410000+0.006654049337+0.003158047686.
\]
Those three terms are observation-attenuation upper bounds for the CSET, TIME, and CONG tiers.  The materializer recomputes them from raw tier budgets and an activation-probability upper bound.  It does not materialize a reference kernel, a cover family, a support floor, a zero-off-support witness, or an MC-EQ transcript projection.

Earlier revisions nevertheless attached fields named \texttt{mc\_eq\_session\_T}, \texttt{per\_step\_epsilon\_t}, and \texttt{conditional\_independence\_assumption} to that row and cited the resulting scalar as a concrete MC-EQ packet.  Those fields did not participate in the arithmetic.  This revision removes them, renames the support pointer to observation attenuation, and adds a validator guard preventing their reintroduction.  The maintained row now states \nolinkurl{semantic_scope=illustrative_tiered_observation_attenuation_not_mceq}, \nolinkurl{composition_evidence_status=assumption_only_no_joint_channel_witness}, and \nolinkurl{publication_eligible=false}.  The $0.056345507023$ value survives only as the tiered observation-attenuation summary that it actually is.

\section{What a real MC-EQ card must contain}
A future deployment card is theorem-bearing only if it publishes all of the following:
\begin{enumerate}[leftmargin=1.5em]
\item the destination family and complete observer transcript projection, including stop/abort outcomes;
\item the initial-law relation to the public reference process;
\item the reference-kernel or cover-family digest;
\item per-history or uniformly valid TV slacks and their evidence source;
\item an explicit mode label: \texttt{tv-only} or \texttt{support-disciplined-ratio};
\item for ratio mode, a zero-off-support-mass witness and a positive $q_t$ floor;
\item the fixed or padded horizon and the exact composition rule;
\item excluded channels and a fail-closed successor trigger.
\end{enumerate}
A card missing item 5 or 6 may still support the transcript-TV theorem.  It may not advertise finite $D_\infty$, pure DP, posterior-mass/PML, or maximal-leakage numbers derived from the reference ratio.

\section{Conclusion}
MC-EQ is salvageable, but only after splitting its claims.  Local TV equalization is a valid and useful transcript-distance contract.  Reference-ratio privacy is a stronger contract that requires support containment and a mass floor.  The off-support tag attack shows why that extra evidence cannot be inferred from a small TV slack, and the worked-example audit shows why decorative fields cannot substitute for a derivation.

\section{Takeaways}
\begin{itemize}[leftmargin=*]
\item Small local TV does not imply finite reference max-divergence.
\item Pairwise TV can stay small while many-destination maximal leakage grows with the destination alphabet.
\item Support containment plus a positive floor converts TV into a valid reference-ratio and maximal-leakage bound.
\item The maintained $0.056345507023$ fallback vector is observation-attenuation evidence, not MC-EQ evidence.
\item No MC-EQ deployment card is publication-ready in this revision.
\end{itemize}

\begin{thebibliography}{9}\small
\bibitem{issa_maxleak}
I.~Issa, A.~B. Wagner, and S.~Kamath.
\newblock An operational approach to information leakage.
\newblock \emph{IEEE Transactions on Information Theory}, 66(3):1625--1657, 2020; arXiv:1807.07878.

\bibitem{levin_peres_wilmer_mixing}
D.~A. Levin and Y.~Peres.
\newblock \emph{Markov Chains and Mixing Times}, second edition.
\newblock American Mathematical Society, 2017.
\end{thebibliography}

\end{document}
